% ============================================================================
% 03_modelos_casos_uso.tex — Cuerpo del entregable "Modelos de Casos de Uso"
% Formato "fully dressed" (Larman, UML y Patrones) para cada caso de uso.
% ============================================================================
\section{Modelos de Casos de Uso}

\subsection{Introducción}

Esta sección desarrolla, en formato completamente vestido (\textit{fully
dressed}) según Craig Larman, cada uno de los doce casos de uso
identificados en el Entregable de Diagrama de Casos de Uso. Los escenarios
se redactan a nivel esencial (intención del actor), sin describir botones,
menús ni otros detalles de la interfaz concreta, que se documentan en el
entregable de mockups. CU-09, CU-10, CU-11 y CU-12 corresponden a los
requisitos RF-03, RF-06, RF-04 y RNF-03 respectivamente.

\subsection{Actores}

\begin{table}[H]
    \centering
    \caption{Actores de LSA Salud}
    \label{tab:actores-mcu}
    \begin{tabularx}{\textwidth}{l l X}
        \toprule
        \textbf{ID} & \textbf{Actor} & \textbf{Descripción} \\
        \midrule
        ACT-01 & Paciente & Persona sorda o hipoacúsica que se comunica mediante LSA. \\
        ACT-02 & Personal de Recepción o Enfermería & Personal de recepción o enfermería que responde al paciente. \\
        ACT-03 & Desarrollador de IA / Equipo de Proyecto & Personal técnico que amplía el vocabulario y reentrena el modelo. \\
        \bottomrule
    \end{tabularx}
\end{table}

% ============================================================================
\subsection{CU-01: Comunicarse mediante Lengua de Señas Argentina}

\begin{longtable}{@{}>{\raggedright\arraybackslash}p{4.1cm} >{\raggedright\arraybackslash}p{\dimexpr\textwidth-4.1cm-2\tabcolsep\relax}@{}}
\caption{CU-01 --- Comunicarse mediante Lengua de Señas Argentina}\\
\toprule
\textbf{Nombre} & Comunicarse mediante Lengua de Señas Argentina \\
\textbf{ID} & CU-01 \\
\textbf{Actor principal} & Paciente \\
\textbf{Personal involucrado e intereses} &
Paciente: quiere expresar una necesidad médica sin depender de un
intérprete humano. \newline
Personal de Recepción o Enfermería: necesita recibir el mensaje del
paciente en español, de forma clara y oportuna. \\
\textbf{Precondiciones} &
El paciente otorgó el consentimiento explícito de cámara (CU-09).
Estación del Paciente abierta, cámara autorizada, el sistema está
operativo y el modelo de reconocimiento está disponible. \\
\textbf{Garantías de éxito} &
El mensaje traducido queda registrado en la conversación compartida y, si
corresponde, se reprodujo en voz. \\
\textbf{Escenario principal de éxito} &
\begin{enumerate}[leftmargin=*, itemsep=2pt]
    \item El paciente se ubica frente a la cámara y comienza a hacer señas.
    \item El sistema captura de forma continua los movimientos del paciente.
    \item El sistema reconoce cada seña realizada y la agrega a la
    secuencia en curso cuando la reconoce con suficiente certeza.
    \item El paciente hace una pausa (baja las manos) al terminar de
    expresar su idea.
    \item El sistema detecta la pausa y da por cerrada la secuencia de
    señas.
    \item El sistema traduce la secuencia a una oración en español natural
    (CU-02: Traducir Secuencia de Señas a Oración Natural).
    \item El sistema sintetiza la oración en voz (CU-03: Sintetizar
    Oración en Voz), si la síntesis está disponible.
    \item El sistema agrega el mensaje traducido a la conversación,
    visible para el personal de recepción o enfermería.
    \item El sistema registra un evento de auditoría anónimo de la
    operación (CU-11: Registrar Evento de Auditoría).
\end{enumerate} \\
\textbf{Extensiones} &
\begin{itemize}[leftmargin=*, itemsep=2pt]
    \item[3a.] La seña reconocida no alcanza el nivel de certeza mínimo: el
    sistema la descarta y continúa observando.
    \item[3b.] La seña reconocida es el estado de reposo (``idle''): el
    sistema no la agrega a la secuencia.
    \item[6a.] El sistema no logra traducir la secuencia en una oración
    clara: descarta el resultado y no genera mensaje.
    \item[7a.] La síntesis de voz no está disponible: el sistema continúa
    sin reproducir audio.
    \item[*a.] En cualquier momento, se pierde la conexión con el servicio
    de reconocimiento: el sistema la restablece automáticamente y reanuda
    el reconocimiento sin intervención del paciente.
\end{itemize} \\
\textbf{Requisitos especiales} & RNF-T-01, RNF-T-05, RNF-01.1 (ver ERS, sección 3.5). \\
\textbf{Variantes de tecnología y datos} &
Captura de landmarks de manos, pose y rostro (blendshapes) mediante
MediaPipe en cada frame; los blendshapes faciales aportan los marcadores no
manuales gramaticales de la LSA (preguntas, negación, intensidad) como
parte de las características de reconocimiento; transporte de los frames
al backend vía WebSocket, con reconexión automática cada 2 segundos ante
una pérdida de conexión. \\
\textbf{Frecuencia} &
Continua durante toda la sesión de atención; se repite cada vez que el
paciente forma una idea completa. \\
\textbf{Temas abiertos} &
El reconocimiento está limitado al vocabulario clínico LSA del dataset del
sistema, que se amplía de forma continua mediante CU-06/CU-07/CU-08 (Modo
Captura); no hay reconocimiento continuo de LSA como lengua completa. \\
\bottomrule
\end{longtable}

% ============================================================================
\subsection{CU-02: Traducir Secuencia de Señas a Oración Natural}

\begin{longtable}{@{}>{\raggedright\arraybackslash}p{4.1cm} >{\raggedright\arraybackslash}p{\dimexpr\textwidth-4.1cm-2\tabcolsep\relax}@{}}
\caption{CU-02 --- Traducir Secuencia de Señas a Oración Natural}\\
\toprule
\textbf{Nombre} & Traducir Secuencia de Señas a Oración Natural \\
\textbf{ID} & CU-02 \\
\textbf{Actor principal} & Paciente (caso de uso incluido desde CU-01) \\
\textbf{Personal involucrado e intereses} &
Paciente: quiere que su idea se entienda correctamente, sin agregados
inventados. \newline
Personal de Recepción o Enfermería: necesita una oración clara y fiel a lo
expresado. \\
\textbf{Precondiciones} &
Existe una secuencia de señas cerrada, con al menos una palabra
reconocida. \\
\textbf{Garantías de éxito} &
Se obtuvo una oración en español, o se determinó que la secuencia no
forma una idea clara. \\
\textbf{Escenario principal de éxito} &
\begin{enumerate}[leftmargin=*, itemsep=2pt]
    \item El sistema recibe la secuencia de palabras sueltas.
    \item El sistema normaliza la secuencia (descarta repeticiones y la
    palabra de reposo).
    \item El sistema busca si la combinación coincide con una plantilla
    médica conocida.
    \item El sistema toma la frase de la plantilla como traducción.
    \item El sistema registra la traducción obtenida, para reutilizarla si
    la misma combinación vuelve a aparecer.
\end{enumerate} \\
\textbf{Extensiones} &
\begin{itemize}[leftmargin=*, itemsep=2pt]
    \item[3a.] No hay coincidencia con ninguna plantilla:
    \begin{itemize}[leftmargin=*, itemsep=2pt]
        \item[3a1.] El sistema busca la combinación entre las traducciones
        generadas recientemente.
        \item[3a2.] Si tampoco encuentra coincidencia, delega la
        construcción de la oración a un modelo de lenguaje local,
        indicándole que respete la persona gramatical y que no agregue
        síntomas no mencionados.
        \item[3a2a.] El modelo de lenguaje no responde a tiempo o falla:
        el sistema usa como traducción las palabras originales, sin
        interrumpir el flujo.
        \item[3a2b.] El modelo de lenguaje determina que la combinación no
        forma una idea clara: el sistema descarta la traducción y no
        genera mensaje.
    \end{itemize}
\end{itemize} \\
\textbf{Requisitos especiales} & RNF-T-01, RF-01.5. \\
\textbf{Variantes de tecnología y datos} &
Plantillas médicas precalculadas (diccionario en el backend); caché de
traducciones recientes; modelo de lenguaje local Qwen2 1.5B servido por
Ollama en \texttt{localhost:11434}; la restricción de no inventar
síntomas se indica mediante instrucciones (\textit{prompt}) al modelo de
lenguaje. \\
\textbf{Frecuencia} & Una vez por cada secuencia de señas cerrada. \\
\textbf{Temas abiertos} &
La cobertura de plantillas médicas sigue el vocabulario clínico del
dataset del sistema. La restricción de no inventar síntomas se refuerza
con una verificación de fidelidad clínica sobre la oración generada
(RF-01.8) antes de mostrarla. \\
\bottomrule
\end{longtable}

% ============================================================================
\subsection{CU-03: Sintetizar Oración en Voz}

\begin{longtable}{@{}>{\raggedright\arraybackslash}p{4.1cm} >{\raggedright\arraybackslash}p{\dimexpr\textwidth-4.1cm-2\tabcolsep\relax}@{}}
\caption{CU-03 --- Sintetizar Oración en Voz}\\
\toprule
\textbf{Nombre} & Sintetizar Oración en Voz \\
\textbf{ID} & CU-03 \\
\textbf{Actor principal} & Paciente (caso de uso incluido desde CU-01) \\
\textbf{Personal involucrado e intereses} &
Personal de Recepción o Enfermería: quiere poder escuchar el mensaje, no
solo leerlo, para atender más rápido. \\
\textbf{Precondiciones} &
Existe una oración traducida; el motor de síntesis de voz está cargado. \\
\textbf{Garantías de éxito} &
Se generó el audio de la oración, o se determinó que la síntesis no está
disponible. \\
\textbf{Escenario principal de éxito} &
\begin{enumerate}[leftmargin=*, itemsep=2pt]
    \item El sistema recibe la oración traducida.
    \item El sistema verifica si ya generó audio antes para esa misma
    oración.
    \item El sistema sintetiza el audio con voz en español rioplatense.
    \item El sistema entrega el audio para su reproducción.
\end{enumerate} \\
\textbf{Extensiones} &
\begin{itemize}[leftmargin=*, itemsep=2pt]
    \item[2a.] El audio ya fue generado antes para la misma oración: el
    sistema reutiliza el audio ya generado en lugar de sintetizarlo de
    nuevo.
    \item[3a.] El motor de síntesis de voz no está cargado: el sistema
    entrega el mensaje sin audio.
\end{itemize} \\
\textbf{Requisitos especiales} & RNF-T-01, RF-01.6. \\
\textbf{Variantes de tecnología y datos} &
Motor Piper TTS con modelo de voz en español argentino
(\texttt{es\_AR-daniela-high}); caché de audio en memoria (hasta 50
oraciones); la interfaz exhibe el texto en pantalla durante una breve
ventana antes de reproducir el audio. \\
\textbf{Frecuencia} & Una vez por cada oración traducida exitosamente. \\
\textbf{Temas abiertos} & Ninguno relevado. \\
\bottomrule
\end{longtable}

% ============================================================================
\subsection{CU-04: Responder al Paciente}

\begin{longtable}{@{}>{\raggedright\arraybackslash}p{4.1cm} >{\raggedright\arraybackslash}p{\dimexpr\textwidth-4.1cm-2\tabcolsep\relax}@{}}
\caption{CU-04 --- Responder al Paciente}\\
\toprule
\textbf{Nombre} & Responder al Paciente \\
\textbf{ID} & CU-04 \\
\textbf{Actor principal} & Personal de Recepción o Enfermería \\
\textbf{Personal involucrado e intereses} &
Personal de Recepción o Enfermería: quiere responder rápido, con el método
que le resulte más cómodo. \newline
Paciente: necesita leer la respuesta en su propia pantalla sin demora. \\
\textbf{Precondiciones} &
Estación de Recepción abierta; existe una conversación activa (puede
estar vacía). \\
\textbf{Garantías de éxito} &
La respuesta del personal de recepción o enfermería queda registrada en la
conversación y visible en la Estación del Paciente. \\
\textbf{Escenario principal de éxito} &
\begin{enumerate}[leftmargin=*, itemsep=2pt]
    \item El personal redacta una respuesta en su propio idioma.
    \item El sistema valida que la respuesta no esté vacía.
    \item El personal confirma el envío.
    \item El sistema agrega la respuesta a la conversación compartida.
    \item El sistema muestra la respuesta en ambas estaciones
    sincronizadas.
    \item El sistema registra un evento de auditoría anónimo de la
    operación (CU-11: Registrar Evento de Auditoría).
\end{enumerate} \\
\textbf{Extensiones} &
\begin{itemize}[leftmargin=*, itemsep=2pt]
    \item[1a.] El personal prefiere dictar la respuesta por voz en
    lugar de escribirla: continúa en CU-05 (Dictar Respuesta por Voz),
    que extiende este caso de uso en el punto de extensión ``ingreso de
    texto''.
    \item[5a.] El paciente o el personal marcan la respuesta como
    incorrecta o ambigua: continúa en CU-10 (Reportar Traducción
    Incorrecta), que extiende este caso de uso en el punto de extensión
    ``revisión de mensaje''.
\end{itemize} \\
\textbf{Requisitos especiales} & RNF-T-05 (sincronización sin demora perceptible). \\
\textbf{Variantes de tecnología y datos} &
Ingreso de texto manual; sincronización entre estaciones mediante
BroadcastChannel (canal \texttt{nexora-chat}, nombre interno heredado del
repositorio de código). \\
\textbf{Frecuencia} &
Cada vez que el personal necesita responder al paciente durante la
atención. \\
\textbf{Temas abiertos} &
Sin historial persistente entre sesiones: si se cierra la pestaña, se
pierde la conversación. \\
\bottomrule
\end{longtable}

% ============================================================================
\subsection{CU-05: Dictar Respuesta por Voz}

\begin{longtable}{@{}>{\raggedright\arraybackslash}p{4.1cm} >{\raggedright\arraybackslash}p{\dimexpr\textwidth-4.1cm-2\tabcolsep\relax}@{}}
\caption{CU-05 --- Dictar Respuesta por Voz}\\
\toprule
\textbf{Nombre} & Dictar Respuesta por Voz \\
\textbf{ID} & CU-05 \\
\textbf{Actor principal} & Personal de Recepción o Enfermería \\
\textbf{Personal involucrado e intereses} &
Personal de Recepción o Enfermería: quiere responder sin escribir, por
ejemplo mientras atiende otras tareas. \\
\textbf{Precondiciones} &
El sistema tiene disponible el reconocimiento de voz; el personal está en
el punto de extensión ``ingreso de texto'' de CU-04. \\
\textbf{Garantías de éxito} &
El texto dictado queda disponible como respuesta, lista para enviarse
(el flujo continúa en CU-04 desde el paso de confirmación). \\
\textbf{Escenario principal de éxito} &
\begin{enumerate}[leftmargin=*, itemsep=2pt]
    \item El personal activa el dictado por voz.
    \item El sistema escucha el habla del personal de forma
    continua.
    \item El sistema transcribe el habla a texto en español rioplatense,
    mostrando resultados parciales a medida que reconoce.
    \item El personal detiene el dictado cuando termina de hablar.
    \item El sistema conserva el texto transcripto, sumándolo a cualquier
    texto que el personal ya hubiera escrito antes de dictar.
    \item El flujo continúa en el paso 2 de CU-04 (confirmar y enviar).
\end{enumerate} \\
\textbf{Extensiones} &
\begin{itemize}[leftmargin=*, itemsep=2pt]
    \item[1a.] El reconocimiento de voz no está disponible: el sistema
    informa la limitación y el personal continúa con el flujo
    principal de CU-04 (escribir).
    \item[2a.] Ocurre un error de micrófono o de reconocimiento: el
    sistema informa el error y detiene la escucha.
\end{itemize} \\
\textbf{Requisitos especiales} & RF-02.2; restricción de compatibilidad de navegador (ver ERS, secciones 3.4 y 3.6). \\
\textbf{Variantes de tecnología y datos} &
Web Speech API nativa del navegador, configurada en español argentino
(\texttt{es-AR}); la Estación de Recepción usa Google Chrome, que la
soporta. \\
\textbf{Frecuencia} & Opcional, a criterio del personal de recepción o enfermería, en cualquier respuesta. \\
\textbf{Temas abiertos} &
La disponibilidad del dictado por voz depende del entorno utilizado por
el personal. \\
\bottomrule
\end{longtable}

% ============================================================================
\subsection{CU-06: Capturar Muestra de Seña para Entrenamiento}

\begin{longtable}{@{}>{\raggedright\arraybackslash}p{4.1cm} >{\raggedright\arraybackslash}p{\dimexpr\textwidth-4.1cm-2\tabcolsep\relax}@{}}
\caption{CU-06 --- Capturar Muestra de Seña para Entrenamiento}\\
\toprule
\textbf{Nombre} & Capturar Muestra de Seña para Entrenamiento \\
\textbf{ID} & CU-06 \\
\textbf{Actor principal} & Desarrollador de IA / Equipo de Proyecto \\
\textbf{Personal involucrado e intereses} &
Desarrollador de IA / Equipo de Proyecto: quiere ampliar el vocabulario
reconocible de forma simple. \newline
Paciente: se beneficia indirectamente de un vocabulario más amplio en el
futuro. \\
\textbf{Precondiciones} &
Estación del Paciente abierta en Modo Captura; cámara autorizada; el
sistema está operativo. \\
\textbf{Garantías de éxito} &
Una o más muestras quedaron grabadas y asociadas a la etiqueta de seña
correspondiente. \\
\textbf{Escenario principal de éxito} &
\begin{enumerate}[leftmargin=*, itemsep=2pt]
    \item El desarrollador entra al Modo Captura.
    \item El desarrollador selecciona o crea la etiqueta de la seña a
    grabar (CU-07: Gestionar Vocabulario de Señas).
    \item El desarrollador indica cuántas repeticiones va a grabar de
    corrido.
    \item El desarrollador inicia la secuencia de grabación.
    \item El sistema anuncia una cuenta regresiva antes de cada
    repetición.
    \item El sistema graba los movimientos del desarrollador durante una
    ventana de tiempo fija.
    \item El sistema guarda la muestra grabada asociada a la etiqueta.
    \item El sistema repite los pasos 5 a 7 hasta completar las
    repeticiones indicadas.
    \item El sistema confirma que las muestras quedaron guardadas y
    actualiza el conteo de muestras por etiqueta.
\end{enumerate} \\
\textbf{Extensiones} &
\begin{itemize}[leftmargin=*, itemsep=2pt]
    \item[2a.] La etiqueta no existe todavía: el sistema la agrega al
    vocabulario disponible en el momento de guardar la primera muestra.
    \item[6a.] El desarrollador detiene la grabación antes de tiempo: el
    sistema cancela esa repetición y no la guarda.
\end{itemize} \\
\textbf{Requisitos especiales} & RF-T-03, RF-T-04. \\
\textbf{Variantes de tecnología y datos} &
Captura mediante MediaPipe: se guardan los landmarks de manos y pose junto
con los 52 coeficientes de expresión facial (blendshapes) de cada frame;
persistencia como archivos \texttt{.npy} en un único árbol de carpetas
(\texttt{dataset/custom/\textless etiqueta\textgreater}). \\
\textbf{Frecuencia} &
Bajo demanda, cuando se necesita ampliar o reforzar el vocabulario. \\
\textbf{Temas abiertos} &
Sin validación automática de la calidad de la seña grabada: la calidad
depende del criterio del desarrollador. \\
\bottomrule
\end{longtable}

% ============================================================================
\subsection{CU-07: Gestionar Vocabulario de Señas}

\begin{longtable}{@{}>{\raggedright\arraybackslash}p{4.1cm} >{\raggedright\arraybackslash}p{\dimexpr\textwidth-4.1cm-2\tabcolsep\relax}@{}}
\caption{CU-07 --- Gestionar Vocabulario de Señas}\\
\toprule
\textbf{Nombre} & Gestionar Vocabulario de Señas \\
\textbf{ID} & CU-07 \\
\textbf{Actor principal} & Desarrollador de IA / Equipo de Proyecto (caso de uso incluido desde CU-06) \\
\textbf{Personal involucrado e intereses} &
Desarrollador de IA / Equipo de Proyecto: necesita saber qué señas ya
existen y cuántas muestras tiene cada una, para decidir qué grabar. \\
\textbf{Precondiciones} & El sistema está operativo. \\
\textbf{Garantías de éxito} &
El desarrollador cuenta con la lista actualizada de etiquetas conocidas y
con la cantidad de muestras disponibles por etiqueta. \\
\textbf{Escenario principal de éxito} &
\begin{enumerate}[leftmargin=*, itemsep=2pt]
    \item El sistema consulta las clases que el modelo reconoce.
    \item El sistema consulta cuántas muestras existen por etiqueta en el
    conjunto de datos activo.
    \item El sistema combina ambas fuentes en una única lista de
    etiquetas conocidas.
    \item El desarrollador elige una etiqueta existente, o crea una
    nueva.
\end{enumerate} \\
\textbf{Extensiones} &
\begin{itemize}[leftmargin=*, itemsep=2pt]
    \item[4a.] El desarrollador crea una etiqueta nueva: el sistema la
    valida (solo letras minúsculas, números y guion bajo, hasta 32
    caracteres) antes de aceptarla.
\end{itemize} \\
\textbf{Requisitos especiales} & RF-T-01, RF-T-02. \\
\textbf{Variantes de tecnología y datos} &
Las estadísticas de muestras se calculan sobre el árbol de datos del
dataset del sistema. \\
\textbf{Frecuencia} &
Cada vez que se abre el Modo Captura, y cada vez que se agrega una
etiqueta. \\
\textbf{Temas abiertos} & Ninguno relevado. \\
\bottomrule
\end{longtable}

% ============================================================================
\subsection{CU-08: Entrenar el Modelo de Reconocimiento}

\begin{longtable}{@{}>{\raggedright\arraybackslash}p{4.1cm} >{\raggedright\arraybackslash}p{\dimexpr\textwidth-4.1cm-2\tabcolsep\relax}@{}}
\caption{CU-08 --- Entrenar el Modelo de Reconocimiento}\\
\toprule
\textbf{Nombre} & Entrenar el Modelo de Reconocimiento \\
\textbf{ID} & CU-08 \\
\textbf{Actor principal} & Desarrollador de IA / Equipo de Proyecto \\
\textbf{Personal involucrado e intereses} &
Desarrollador de IA / Equipo de Proyecto: quiere que el modelo aprenda las
señas nuevas capturadas. \newline
Paciente y Personal de Recepción o Enfermería: se benefician de un modelo
más preciso y con más vocabulario. \\
\textbf{Precondiciones} &
Existen muestras de entrenamiento capturadas (CU-06) para al menos una
clase de seña. \\
\textbf{Garantías de éxito} &
Se generó una nueva versión del modelo de reconocimiento, entrenada con
el conjunto de datos disponible, lista para usarse en las próximas
sesiones. \\
\textbf{Escenario principal de éxito} &
\begin{enumerate}[leftmargin=*, itemsep=2pt]
    \item El desarrollador inicia el proceso de entrenamiento fuera de
    línea.
    \item El sistema releva las clases de seña disponibles en el conjunto
    de datos.
    \item El sistema separa las muestras en un conjunto de entrenamiento
    y uno de validación.
    \item El sistema entrena el modelo de reconocimiento durante un
    número fijo de rondas de aprendizaje.
    \item El sistema evalúa la exactitud del modelo sobre el conjunto de
    validación en cada ronda.
    \item El sistema conserva la versión del modelo con mejor exactitud
    de validación.
    \item El sistema deja disponible el modelo entrenado y su lista de
    clases para cargarlos en el próximo inicio.
\end{enumerate} \\
\textbf{Extensiones} &
\begin{itemize}[leftmargin=*, itemsep=2pt]
    \item[2a.] No se encuentran clases ni muestras en el conjunto de
    datos: el sistema informa la situación y no entrena.
\end{itemize} \\
\textbf{Requisitos especiales} & RF-T-05. \\
\textbf{Variantes de tecnología y datos} &
Aumento de datos por recorte aleatorio, ruido gaussiano y escalado
durante el entrenamiento; separación 80/20 con semilla fija para
reproducibilidad; se ejecuta desde la línea de comandos del servidor, sin
interfaz web ni disparo automático. \\
\textbf{Frecuencia} &
Bajo demanda del desarrollador, típicamente después de acumular nuevas
muestras de varias sesiones de captura. \\
\textbf{Temas abiertos} &
El entrenamiento se inicia manualmente, sin disparo automático. \\
\bottomrule
\end{longtable}

% ============================================================================
\subsection{CU-09: Otorgar Consentimiento Explícito de Cámara}
\textit{Cubre el requisito RF-03.}

\begin{longtable}{@{}>{\raggedright\arraybackslash}p{4.1cm} >{\raggedright\arraybackslash}p{\dimexpr\textwidth-4.1cm-2\tabcolsep\relax}@{}}
\caption{CU-09 --- Otorgar Consentimiento Explícito de Cámara}\\
\toprule
\textbf{Nombre} & Otorgar Consentimiento Explícito de Cámara \\
\textbf{ID} & CU-09 \\
\textbf{Actor principal} & Paciente \\
\textbf{Personal involucrado e intereses} &
Paciente: necesita entender, antes de que se lo grabe, para qué se usa la
cámara y poder negarse. \newline
Sanatorio (Comité de Ética): necesita que ningún paciente sea grabado sin
autorización explícita. \\
\textbf{Precondiciones} &
El paciente abrió la Estación del Paciente; la cámara todavía no fue
activada. \\
\textbf{Garantías de éxito} &
La cámara se activa únicamente si el paciente otorgó el consentimiento; si
lo rechaza, la cámara permanece inactiva y CU-01 no puede comenzar. \\
\textbf{Escenario principal de éxito} &
\begin{enumerate}[leftmargin=*, itemsep=2pt]
    \item Al abrir la Estación del Paciente, el sistema presenta al
    paciente una solicitud de consentimiento, antes de solicitar acceso a
    la cámara.
    \item El sistema explica, en lenguaje simple, qué captura la cámara,
    para qué se usa y que no se guarda video ni imágenes (RNF-03).
    \item El paciente otorga el consentimiento.
    \item El sistema registra el consentimiento como evento de auditoría
    anónimo, sin datos identificatorios, y habilita la captura de cámara.
    \item El flujo continúa en el paso 1 de CU-01.
\end{enumerate} \\
\textbf{Extensiones} &
\begin{itemize}[leftmargin=*, itemsep=2pt]
    \item[3a.] El paciente rechaza el consentimiento: el sistema registra
    el evento de auditoría correspondiente, no activa la cámara y ofrece
    al paciente la opción de reconsiderar su decisión; CU-01 no puede
    iniciar.
    \item[3b.] El acceso a la cámara es denegado después de un
    consentimiento otorgado: el sistema informa el error sin activar el
    flujo de reconocimiento.
\end{itemize} \\
\textbf{Requisitos especiales} & RF-03.1; RNF-03.1 (sin almacenamiento de video). \\
\textbf{Variantes de tecnología y datos} &
Diálogo modal de consentimiento (Pantalla P-05) dentro de
\texttt{PacientePage} (\texttt{web/app/paciente/page.tsx}), mostrado
sobre el resto de la página al entrar a \texttt{/paciente}, con los
botones ``Acepto y continúo'' / ``No acepto''; el estado
\texttt{consentState} (\texttt{pending}/\texttt{accepted}/\texttt{declined})
se pasa como \texttt{enabled} al hook de captura (\texttt{useMediaPipe}),
que solo invoca \texttt{navigator.mediaDevices.getUserMedia()} cuando
\texttt{enabled === true}; el evento de auditoría registra
\texttt{type: "consent"} con \texttt{consent: "accepted"} o
\texttt{"declined"}; un permiso de cámara denegado se refleja como
\texttt{status: "error"} en el hook de captura. \\
\textbf{Frecuencia} & Una vez por sesión de la Estación del Paciente. \\
\textbf{Temas abiertos} &
El consentimiento no se conserva entre recargas de la página, en línea
con la política de no almacenamiento persistente (RNF-03); el texto legal
de la solicitud se valida con el Comité de Ética del Sanatorio durante la
instalación de cada estación. \\
\bottomrule
\end{longtable}

% ============================================================================
\subsection{CU-10: Reportar Traducción Incorrecta o Ambigua}
\textit{Cubre el requisito RF-06. Extiende a CU-01 y CU-04.}

\begin{longtable}{@{}>{\raggedright\arraybackslash}p{4.1cm} >{\raggedright\arraybackslash}p{\dimexpr\textwidth-4.1cm-2\tabcolsep\relax}@{}}
\caption{CU-10 --- Reportar Traducción Incorrecta o Ambigua}\\
\toprule
\textbf{Nombre} & Reportar Traducción Incorrecta o Ambigua \\
\textbf{ID} & CU-10 \\
\textbf{Actor principal} & Paciente, Personal de Recepción o Enfermería \\
\textbf{Personal involucrado e intereses} &
Paciente y Personal: quieren señalar cuando un mensaje traducido no
representa lo que se quiso decir, evitando que el personal tome decisiones
clínicas o administrativas sobre una traducción errónea, y que el equipo
de desarrollo mejore el modelo. \newline
Desarrollador de IA / Equipo de Proyecto: necesita casos reales
identificados para priorizar el reentrenamiento del modelo (CU-08). \\
\textbf{Precondiciones} &
Existe al menos un mensaje en la conversación (generado por CU-01 o
CU-04). \\
\textbf{Garantías de éxito} &
El reporte queda registrado asociado al mensaje señalado, el mensaje queda
marcado como traducción reportada en ambas caras de la estación, y se
instruye al paciente a repetir la seña; todo ello sin almacenar datos
identificatorios del paciente. \\
\textbf{Escenario principal de éxito} &
\begin{enumerate}[leftmargin=*, itemsep=2pt]
    \item El actor revisa un mensaje del paciente en la conversación
    (punto de extensión ``revisión de mensaje'' de CU-01 o CU-04).
    \item El actor indica que desea reportar la traducción del mensaje.
    \item El sistema envía el reporte con el identificador del mensaje, su
    texto, la estación de origen y, si están disponibles, las palabras
    crudas reconocidas.
    \item El sistema valida la información y registra el reporte asociado
    al mensaje señalado.
    \item El sistema notifica el reporte a ambas caras de la estación.
    \item El sistema actualiza el mensaje en las dos caras: lo marca como
    traducción reportada y exhibe un aviso operativo (al paciente
    se le solicita repetir la seña y al personal se le advierte no actuar
    sobre esa traducción).
    \item El flujo vuelve al punto de extensión de origen (CU-01 o
    CU-04).
\end{enumerate} \\
\textbf{Extensiones} &
\begin{itemize}[leftmargin=*, itemsep=2pt]
    \item[3a.] El envío del reporte falla (red o validación): el sistema
    informa que no pudo enviarse y permite reintentar sin alterar el estado
    del mensaje.
\end{itemize} \\
\textbf{Requisitos especiales} & RF-06.1; RNF-03.1 (reporte anónimo, sin almacenar la conversación completa). \\
\textbf{Variantes de tecnología y datos} &
Botón ``Reportar traducción'' en \texttt{ChatBubble}
(\texttt{web/components/chat/ChatBubble.tsx}), visible en mensajes del
paciente; al reportar invoca \texttt{POST /report} y luego sincroniza el
evento \texttt{REPORTED} mediante \texttt{useChatChannel}
(\texttt{BroadcastChannel("nexora-chat")}). El backend valida el payload
(\texttt{validate\_report\_payload}) y delega la persistencia al módulo
\texttt{storage.py} sobre la base de datos SQLite (\texttt{data/lsa\_salud.db}).
El mensaje pasa a estilo tachado con fondo de advertencia y textos
contextualizados por estación. \\
\textbf{Frecuencia} &
Bajo demanda, cada vez que el paciente o el personal detectan una
traducción incorrecta. \\
\textbf{Temas abiertos} &
Ninguno. El reporte captura el texto del mensaje, su identificador, la
estación y las palabras reconocidas como insumo directo para priorizar el
reentrenamiento (CU-08). \\
\bottomrule
\end{longtable}

% ============================================================================
\subsection{CU-11: Registrar Evento de Auditoría}
\textit{Cubre el requisito RF-04. Subfunción incluida por CU-01 y CU-04, sin actor humano.}

\begin{longtable}{@{}>{\raggedright\arraybackslash}p{4.1cm} >{\raggedright\arraybackslash}p{\dimexpr\textwidth-4.1cm-2\tabcolsep\relax}@{}}
\caption{CU-11 --- Registrar Evento de Auditoría}\\
\toprule
\textbf{Nombre} & Registrar Evento de Auditoría \\
\textbf{ID} & CU-11 \\
\textbf{Actor principal} & Ninguno --- subfunción interna del sistema (ver \texttt{IDS.md}, Sección~1, sobre la decisión de no modelar un actor de auditoría). \\
\textbf{Personal involucrado e intereses} &
Dirección del Sanatorio y Equipo de Proyecto: necesitan métricas
operativas anónimas para evaluar el impacto del sistema, sin que eso
implique acceso a datos identificatorios del paciente. \\
\textbf{Precondiciones} &
CU-01 o CU-04 se está ejecutando y alcanzó un punto de inclusión de este
caso de uso. \\
\textbf{Garantías de éxito} &
Un evento de auditoría anónimo quedó registrado, con tipo de operación,
marca temporal y resultado, sin datos identificatorios ni imágenes o
video. \\
\textbf{Escenario principal de éxito} &
\begin{enumerate}[leftmargin=*, itemsep=2pt]
    \item El sistema completa una operación relevante: predicción de una
    seña, traducción y síntesis de voz, reporte de una traducción,
    consentimiento de cámara, traducción mostrada en pantalla o fin de
    atención (CU-12).
    \item El sistema arma un evento de auditoría anónimo con el tipo de
    operación, la marca temporal y los datos propios del tipo de evento,
    sin identidad del paciente ni contenido de imagen, video o texto de
    la conversación.
    \item El sistema persiste el evento.
\end{enumerate} \\
\textbf{Extensiones} &
\begin{itemize}[leftmargin=*, itemsep=2pt]
    \item[--] No hay un camino de excepción propio: si la información de
    un evento originado en la Estación del Paciente o de Recepción no
    pasa la validación, el sistema simplemente no lo registra, sin
    interrumpir la operación que lo originó.
\end{itemize} \\
\textbf{Requisitos especiales} & RF-04.1; RNF-03.1 (anonimato, sin imágenes ni video). \\
\textbf{Variantes de tecnología y datos} &
Persistencia relacional embebida: la función \texttt{\_log\_audit} de
\texttt{app.py} delega en \texttt{storage.log\_audit\_event} la inserción en
la tabla \texttt{audit\_event} de la base de datos SQLite
(\texttt{data/lsa\_salud.db}), consistente con el resto del sistema (ver
Entregable 09, Modelo de Datos). Las
operaciones que disparan el registro corresponden a los endpoints
\texttt{/predict}, \texttt{/ws/predict}, \texttt{/process\_sentence},
\texttt{/report} y \texttt{POST /audit} (consentimiento y traducción
mostrada en pantalla); el tipo de evento toma los valores
\texttt{prediction}, \texttt{sentence}, \texttt{report}, \texttt{consent},
\texttt{translation\_displayed} o \texttt{fin\_atencion}, con campos
propios por tipo (p.\ ej.\ \texttt{lstm\_ms} y \texttt{confidence} para
\texttt{prediction}); el envío desde el frontend es \textit{fire-and-forget}
y, si el payload no pasa la validación, \texttt{POST /audit} responde 400
sin registrar el evento; las métricas agregadas (conteos, latencias
p50/p95/máx.) se exponen bajo demanda mediante \texttt{GET /metrics}. \\
\textbf{Frecuencia} & Una vez por cada operación de sistema relevante. \\
\textbf{Temas abiertos} &
El sistema agrega los eventos acumulados en métricas (conteos, latencias
p50/p95/máx.) disponibles bajo demanda; el Equipo de Proyecto las exporta
periódicamente para evaluar el impacto operativo del sistema. \\
\bottomrule
\end{longtable}

% ============================================================================
\subsection{CU-12: Finalizar Atención y Preparar la Estación para el Próximo Paciente}
\textit{Cubre el requisito RNF-03.}

\begin{longtable}{@{}>{\raggedright\arraybackslash}p{4.1cm} >{\raggedright\arraybackslash}p{\dimexpr\textwidth-4.1cm-2\tabcolsep\relax}@{}}
\caption{CU-12 --- Finalizar Atención y Preparar la Estación para el Próximo Paciente}\\
\toprule
\textbf{Nombre} & Finalizar Atención y Preparar la Estación para el Próximo Paciente \\
\textbf{ID} & CU-12 \\
\textbf{Actor principal} & Personal de Recepción o Enfermería \\
\textbf{Personal involucrado e intereses} &
Personal de Recepción o Enfermería: necesita dejar la estación lista para
el próximo paciente sin arrastrar datos de la atención anterior. \newline
Paciente: se beneficia de que ningún dato de su atención quede accesible
para el siguiente paciente (RNF-03). \\
\textbf{Precondiciones} &
Se inició una atención en la estación (CU-01 y/o CU-04, en curso o ya
finalizada). \\
\textbf{Garantías de éxito} &
La conversación queda vacía en las dos caras de la estación y no persiste
ningún dato de la sesión atendida. \\
\textbf{Escenario principal de éxito} &
\begin{enumerate}[leftmargin=*, itemsep=2pt]
    \item El personal de recepción o enfermería indica que finaliza la
    atención.
    \item El sistema borra la conversación mostrada en esa cara de la
    estación.
    \item El sistema notifica el cierre a la otra cara de la estación
    mediante el mismo canal de sincronización que usa CU-04.
    \item La otra cara de la estación borra su propia copia de la
    conversación.
    \item El sistema registra un evento de auditoría anónimo de fin de
    atención (CU-11: Registrar Evento de Auditoría).
    \item La estación queda lista para atender al próximo paciente, sin
    ningún dato de la sesión anterior.
\end{enumerate} \\
\textbf{Extensiones} &
\begin{itemize}[leftmargin=*, itemsep=2pt]
    \item[--] No hay un camino de excepción propio: borrar la
    conversación es una operación local e inmediata en cada cara de la
    estación, sin dependencia de red ni de un servicio externo.
\end{itemize} \\
\textbf{Requisitos especiales} & RNF-03 (privacidad por diseño); RF-04.1 (auditoría). \\
\textbf{Variantes de tecnología y datos} &
Botón ``Finalizar atención'' (Pantalla P-03); mismo canal de
sincronización entre las dos caras de la estación que usa CU-04 para
propagar mensajes. \\
\textbf{Frecuencia} & Una vez por paciente atendido. \\
\textbf{Temas abiertos} & Ninguno relevado. \\
\bottomrule
\end{longtable}
